\section{引言}\label{sec:intro}
交通预测为路网运行管理、出行规划与拥堵预警提供前瞻性信息。然而，预测系统能够获得的历史观测并不总是连续且完整。随着传感网络的建设与维护，路网中的传感器处于不同的观测状态~\citep{zhou2025mogernn}：持续可用传感器（AAS）在训练与推理期间稳定提供读数；新增传感器（NAS）在训练期没有记录，接入后才开始积累读数；失效传感器（FS）原本正常工作，后因设备故障或通信异常中断读数；虚拟传感器（VS）则是没有可读观测、只需要预测的位置。已有关于未观测位置预测和动态传感器网络的研究已认识到这些问题~\citep{mei2023uncertainty,zhou2025mogernn,zhao2026structure}。同一位置在不同预测窗口中可以处于不同状态，所提供的信息条件并不相同。本文将历史窗口内实际可用观测的位置、数量及其时间分布统称为观测支持（OS）。OS 的变化不仅产生需要处理的空缺，也决定模型能从本地历史与关联位置获得怎样的预测依据。

围绕不完整交通观测，已有研究形成了互补的解决途径。缺失感知方法通过图循环重建、条件扩散或联合重建与预测，利用时空关联补充不可用信息~\citep{cini2022filling,liu2023pristi,zuo2023graph,yu2024ginar}；归纳与部分感知预测将预测范围扩展至缺少训练信号或推理期持续观测的位置~\citep{mei2023uncertainty,zhou2025mogernn,liu2025spatiotemporal,zhao2026structure}；自适应图学习及空间分块表示则改善空间关系和多尺度上下文的组织~\citep{yu2018spatio,wu2019graph,bai2020adaptive,jiang2023pdformer,fang2025efficient}。这些研究已经证明，缺失处理、空间信息传递与未观测位置预测能够协同工作。因而，进一步的关键并非判断模型是否能够接受缺失输入，而是明确：当 OS 在预测窗口内变化时，模型应如何相应调整信息利用方式。

这一问题在归纳预测中尤其突出。目标节点没有参与动态信号训练，并不意味着它在推理时始终没有观测；同样，训练期已有记录也不保证推理时始终拥有完整历史。现有方法分别强调训练信号覆盖、缺失恢复、传感器配置适应或空间交互，但这些能力并不直接回答“同一目标对应不同 OS 时，预测计算应如何变化”。即使可用观测数相同，近期集中、早期集中和分散分布的历史也可能对未来状态提供不同依据。本文因此研究\textbf{动态 OS 下的归纳时空预测}：在路网拓扑已知、部分节点动态信号不参与训练的条件下，利用变化的 OS 组织空间推理与预测上下文。该设置连接节点归纳与窗口内 OS 的变化，强调固定模型参数下的预测行为。

我们的核心认识是，\textbf{OS 应同时作为输入有效性的描述和预测计算的条件}。观测掩码标明哪些读数可用；在此基础上，观测角色与状态进度可以为模型组织历史信息提供显式的结构化线索。观测角色记录每个节点在每一步属于 AAS、NAS、VS 还是 FS；状态进度（state progress）描述接入或中断已持续的比例。它们提供一种可检验的归纳偏置，而不要求包含完整掩码之外的额外信息。沿这一思路，历史补全负责提供可用表示，空间传播决定如何交换关联信息，跨节点时间上下文则为预测过程提供条件。三者的协作，使模型能够在本地证据与空间证据之间建立依赖当前 OS 的信息利用方式。

据此，我们提出 ALFNet（Adaptive Lifecycle Forecasting Network），其中“自适应”指计算随输入的 OS 而变化，推理时模型参数保持固定。ALFNet 由三个模块组成，分别对应 OS 变化带来的三个问题。第一，不同节点、不同时刻的缺失程度各不相同，若各分支分别处理缺失，同一窗口会得到不一致的输入。观测支持感知历史补全（SAHC）因此依据可见读数、观测角色、状态进度与路网关联，对全部节点只补全一次历史，供另外两个模块共享。第二，当目标节点的本地历史很少或已经中断时，预测需要来自关联位置的依据。空间上下文时间条件分支（ALFTemporal）把节点组织为空间分区，在分区之间与分区内部交换时间特征，生成逐节点的条件特征，并经特征级线性调制（FiLM）作用于图编码状态；它本身不产生预测。第三，补全之后，真实读数与补全值在数值上难以区分，同一种图传播方式也难以同时适合持续可用、刚接入与刚失效的节点。观测条件化图预测（ALFGraph）因此由每一步的观测角色、状态进度与掩码决定五种图传播算子的组合，在编码与解码中逐步调整空间信息的交换方式，并以残差自回归方式生成预测；输出端再由末步观测角色控制一个有界的图修正。三个模块组合后，SAHC 提供一致的输入，ALFTemporal 提供跨节点的预测依据，ALFGraph 决定这些信息在每个节点、每一步如何参与传播，使同一组固定参数能够随 OS 调整对本地历史与空间上下文的利用。本文的重点是这一条件化组织方式；所用的图循环补全、扩散卷积门控循环、空间分区注意力与特征级线性调制等基础算子均来自已有工作~\citep{cini2022filling,fang2025efficient,perez2018film,li2018diffusion,cho2014learning,zhang2023crossformer}。

本文的贡献概括为三点：（1）围绕拓扑已知、训练信号覆盖不完整的路网，将窗口内变化的 OS 纳入归纳预测设置，按目标节点的观测角色轨迹定义零历史（Unseen）、稳定观测（Steady）、新增接入（Join）与故障中断（Fault）四类场景，并按事件持续步数细分 Join 与 Fault 的阶段，以刻画不同 OS 下的预测要求；（2）提出 ALFNet，将 SAHC、ALFTemporal 与 ALFGraph 组织为以 OS 为条件的协同预测框架；（3）在 METR-LA、PEMS-BAY、PEMS04 与 PEMS08 四个数据集上开展统一的场景与阶段评价，并结合组件消融、配对观测干预和样本层面分析，考察总体收益、机制归因及 OS 对预测行为的影响。

结果表明，在统一的数据暴露与评价协议下，ALFNet 在四个真实世界数据集上跨场景平均误差（LifeAvg）的三项误差指标均优于八个代表性比较方法，平均绝对误差（MAE）降低7.86\%–15.63\%；在本地证据逐步积累或中断的 Join 与 Fault 场景下，改进在四个数据集的三项指标上一致。
